\clearpage
\Question{Big-O Notation}
%% Task 3: comment version A or B in alternate semesters

Recall the definition for big-O:
\begin{quote}\em
  $g \in O(f)$ if there is a real constant $c > 0$ and some natural
  number $n_0$ such that
  for every $n \ge n_0$ we have $g(n) \le c \cdot f(n)$.
\end{quote}

\begin{parts}
\newcommand{\poly}
  {3n^2 + 7n + 4} % [c=4,n0=8; c=5,n0=4; ...] F19 S18 F16 S16
%  {2n^2 + 5n + 1}  % [c=3,n0=6; c=4,n0=3; ...] F18 S17
%  {4n^2 + 2n + 3}  % [c=5,n0=0; c=6,n0=1; ...] S19 F17

\RUBRIC
NOTE: the function changes from semester to semester.  Make sure the
      rubrics match the function.
ENDRUBRIC


\part[2]\TAGS{big-o}
Demonstrate graphically that $\poly \in O(n^2)$ by finding values for
$c$ and $n_0$ that satisfy the definition above, where $g(n) = \poly$
and $f(n) = n^2$. Draw a picture to illustrate that $cn^2$ acts as an
upper bound to $\poly$ for all $n \ge n_0$ using your values for $c$
and $n_0$. Be sure to label the horizontal and vertical axes.

If you choose to insert a plot (or a picture) generated using an
external application, please \textbf{paste it in the space provided} ---
do not insert it as a new page.
\begin{framed}
\bigskip
$c = \uanswer{24.5em}{5}$
\par\bigskip
$n_0 = \uanswer{24em}{15}$

\bigskip
\begin{center}
\begin{tikzpicture}
\draw[step=0.5cm,color=gray] (-1,-1) grid (13,9);
\end{tikzpicture}

\smallskip
\end{center}
\end{framed}

\RUBRIC
Part (a)
TAGS: big-o

Gradescope rubric:
+1pt Reasonable n0 and c (Examples: c=10 and n0=1, ...)
+1pt Clean plot with labeled axes
-0.5pt Curves not labeled

Commentary:
. Answers:
  To check answers: https://www.mathsisfun.com/quadratic-equation-solver.html

  To come up with answers:
  - set c = 1 + quadratic coefficient (or any strictly positive real intead of 1)
  - simplify the inequality cn^2 - (given poly) >= 0
  - plug the coefficients in above web site
  - pick any n0 on the right where the parabola is positive

  Grading: 1 pt for constants, 1 pt for a reasonable picture to
  accompany the constants.  NOTE: If necessary, you can go to
  fooplot.com and plot the two functions if you need to verify.
ENDRUBRIC

\newpage
\part[0\half]\TAGS{big-o}
In one sentence, explain why $\poly \notin O(n)$.
\begin{framed}
\ifprintanswers{\color{\answerColor}
  Since a quadratic function increases more rapidly than a linear
  function, there is no constant that can guarantee that the linear
  function will always be an upper bound to the quadratic function.
}\else~\vspace{1.9in}\fi
\end{framed}

\RUBRIC
Part (b)
TAGS: big-o

Gradescope rubric:
+ 0.25 pt  Sentence says that the first function grows faster than the second
+ 0.25 pt  Sentence says that there is no constant that can guarantee that the linear function will always be an upper bound to the quadratic function

Commentary:
. Since a quadratic function increases more rapidly than
  a linear function, there is no constant that can guarantee that the
  linear function will always be an upper bound to the quadratic
  function.

   Grading: 1 pt for anything that indicates that the linear function
   can't be an upper bound to a quadratic since it doesn't grow as
   fast. 1/2 pt for something in the ball park but says something
   vague or inconsistent (try to avoid this, and discuss to maintain
   consistency among graders). 0 pts for something just plain
   wrong. Students might try to justify using the formal definition,
   which is ok if they state something that there's no c that can be
   found that will make the linear function an upper bound...
ENDRUBRIC

\enlargethispage{5ex}
\part[1]\TAGS{complexity}
Determine the asymptotic complexity of the following function using
big-O notation as a function of its arguments in its simplest, tightest
form.

%\bgroup

\renewcommand{\a}
%{a} % up to S18
{u} % S19
\renewcommand{\b}
%{b} % up to S18
{v} % S19
\newcommand{\x}
%{x} % up to S18
{y} % S19

\begin{lstlisting}
int mystery(int [*\a*], int [*\b*])
//@requires [*\color{\contractColor}\a*] > 0 && [*\color{\contractColor}\b*] > 0;
{
    int [*\x*] = 0;
    for (int i = 1; i < 10; i++) {
        int j = [*\a*];

        while (j > 0) {
            int k = [*\b*];

            while (k > 0) {
                [*\x*] = [*\x*] + i * j * k;
                k = k / 2;
            }

            j--;
        }
    }

    return [*\x*];
}
\end{lstlisting}

\begin{framed}\vspace{0.25in}
$O(\uanswer{16em}{$\a \log \b$})$
\end{framed}

\RUBRIC
Part (mystery A)
TAGS: big-o

Gradescope rubric:
+ 0.5 pts Big-O is in simplest form: (O(u log v), not O(u log_2 v)
+ 0.5 pts Big-O is correct and in tightest form

Commentary:
- check name of variables (varies from semester to semester)
. 1/2 pt if student mixes up a and b or gets part of it right
  (e.g. O(ab)).
. 1/2 pt if student has O(10 * a log b) since 10 doesn't matter
  here, or if they include the base 2 e.g. O(a log_2 b)
ENDRUBRIC

\egroup % S19, before F18
\bgroup
\renewcommand{\a}
{a} % up to F18
\renewcommand{\b}
{b} % up to F18
\newcommand{\A}
{A} % up to F18
\newcommand{\n}
{25} % up to F18

\begin{lstlisting}
int mystery(int[] [*\A*], int [*\a*], int [*\b*])
//@requires 0 <= [*\color{\contractColor}\a*] && [*\color{\contractColor}\a*] < \length([*\color{\contractColor}\A*]);
//@requires 0 <= [*\color{\contractColor}\b*] && [*\color{\contractColor}\b*] < \length([*\color{\contractColor}\A*]);
{
    int r = 0;
    for (int i = 1; i < [*\n*]; i++) {
        int j = [*\a*];

        while (j > 0) {
            int k = [*\b*];

            while (k > 0) {
                r = r + i * [*\A*][j] * [*\A*][k];
                k--;
            }

            j = j / 2;
        }
    }

    return r;
}
\end{lstlisting}

\begin{framed}\vspace{0.25in}
$O(\uanswer{16em}{$b \log a$})$
\end{framed}

\RUBRIC
Part (mystery B)
TAGS: big-o

Gradescope rubric:
+0.5pt Big-O is in simplest form: O(b log a), not O(b log_2 a)
+0.5pt Big-O is correct and in tightest form

Commentary:
Grading: 1 pt for correct answer with no unnecessary constants
. 1/2 pt if student mixes up a and b or gets part of it right
  (e.g. O(ab)).
. 1/2 pt if student has O(25 * b log a) since 25 doesn't matter
  here, or if they include the base 2 e.g. O(b log_2 a)
ENDRUBRIC

\egroup % F19 F18
\end{parts}

\bigskip%
\textbf{Reminder:} as discussed in recitation, in this class when we
are dealing with logarithms, we consider the ``simplest form'' to be
the one \emph{without the base}.  Therefore, we prefer $O(\log n)$
over big-$O$ descriptions like $O(\log_4 n)$ or $O(\ln n)$.
